\chapter{运行载体层：计算介质、资源约束与跨尺度接口}
\label{chp:medium_layer_ch}

智能过程并不在抽象空间中自行运行。无论系统依托神经组织、电子设施、群体协作网络，还是把环境痕迹纳入闭环，它都必须借助某种可更新、可保持、可传递并可读出的载体。载体不会凭材料名称自动获得语义，却会通过时延、容量、精度、可塑性、故障方式与接口条件，决定哪些状态轨迹能够被实现和复现。本章据此把介质问题重建为运行载体问题：我们先规定计算对象、观测尺度和任务判据，再讨论不同实现怎样约束智能状态动力学。HSF--HD在这里提供跨平台的一般状态语言；AgentOS只是其中一个工程实例，并非所有智能系统的唯一实现。所谓智能的一般“空气动力学”，在本章具体指向对稳定域、扰动、资源和失效边界的研究，而不把智能等同于任何物理流体。

\section{计算载体谱系：生物网络、电子系统与环境介质}
\label{sec:physical_substrate_spectrum}

我们把运行载体定义为四元组 $P=(V,C,M,Q)$。其中 $V$ 是执行状态变换的处理单元集合，$C$ 是信息能够通过的有向通道集合，$M$ 是在给定时间尺度上保持状态的存储集合，$Q$ 是感知、行动和观测接口集合。这个定义不要求处理单元必须是神经元或晶体管，也不要求存储位于皮肤或机箱之内；它只要求说明一次更新在哪里发生、依赖何种输入、结果保持多久，以及观察者凭什么判断更新有效。对任务 $\tau$，载体支持的可实现轨迹写为
\[
\mathcal T_{P,\tau}=\{(x_0,u_0,x_1,\ldots,x_T):x_{k+1}=F_P(x_k,u_k;\theta_P)\},
\]
其中 $x_k$ 是离散时刻 $k$ 的状态，$u_k$ 是接口输入，$F_P$ 是受连接、存储和资源限制的更新算子，$\theta_P$ 汇集平台参数。步长必须由实现声明；神经记录的一步、GPU内核的一步和蚁群一次信息素更新并不具有天然相同的持续时间，因此跨载体比较应比较任务轨迹、误差与代价，不能把微观事件逐项等同。

生物神经系统是高并行、局部连接、异步且持续可塑的载体。其计算状态同时受神经活动、突触连接、调质信号与身体反馈影响。电活动和化学活动是实现条件，却不等于语义本身；振荡同步、放电率或网络拓扑只有在特定任务和记录协议下才支持相应经验结论。生物载体通常具有紧密的感知行动闭环、跨时间尺度学习和一定损伤重组能力，但也存在测量分辨率有限、个体差异显著、状态难以复制、训练历史难以控制等限制。脑区分工和节律现象可以提供计算假设，却不能直接规定机器模块与其机制身份相同。

电子计算系统以时钟、指令、缓存、互连和持久化设备组织更新。它便于记录状态、复制实验与隔离变量，但同样受到内存层级、通信拥塞、数值误差、并行一致性和功耗预算约束。大模型的一次前向计算、多代理系统的消息交换和数据库事务虽都由电子器件实现，却属于不同协议；统称为“硅基心智”会遮蔽真正决定轨迹的接口与调度规则。必须区分硬件可用资源、软件可见资源和任务实际占用：峰值吞吐不等于端到端有效吞吐，显存容量也不等于可无损保持的语义上下文。

环境介质构成第三种重要实现。化学扩散、信息素、纸面记号、共享数据库与公共黑板都能承担外部保持和协调功能。蚁群的信息素轨迹不是集中处理器中的向量，却能改变后续个体的转移概率；人的笔记不在大脑内部，却能稳定跨时段目标；AgentOS中的外部记忆 $M_e$ 同样属于扩展实现。只要外部痕迹可可靠写入、保持、寻址并重新影响决策，它就是闭环载体的一部分。系统边界应由因果干预和权限契约确定，而不是凭空间位置决定。该谱系保留生物、电子和化学实现的差异，同时避免宣称它们整体同构；真正可比较的是操作、接口、故障与任务后果。

载体分类还必须允许混合。机器人可用电子计算维护模型，用机械身体获取反馈，用云端数据库保存经验，并借人类协作完成超出自身能力的步骤；人类也会借语言、组织和工具扩展内部处理。我们不因一个部件位于外部就把它排除，也不因它参与闭环就取消主体边界。纳入条件应包括可达性、因果贡献、更新权限、失效影响与恢复责任。通过逐项断开接口并测量性能变化，可以估计外部部件对任务的真实贡献。这样，运行载体层得到的是可检验的系统分解，而不是“万物皆心智”的无限扩张。

\section{载体约束参数：时延、吞吐、保持与探索强度}
\label{sec:physical_constants}

旧有表述把平台差异概括为“认知光速”“介质黏度”和“认知温度”。这些名称有启发性，却容易把通信时延、状态保持、随机探索与焦耳能耗混为一谈。我们改用可测量参数。设载体图为 $G=(V,C)$，通道 $e$ 的时延为 $\ell_e$ 秒，节点 $v$ 的服务时间为 $s_v$ 秒，通道有效吞吐为 $b_e$，其单位可以是比特每秒、令牌每秒或任务消息每秒，但同一次比较必须固定口径。对依赖路径集合 $\Pi_\tau$，无排队近似下的关键路径下界为
\[
L_\tau^{\min}=\max_{\pi\in\Pi_\tau}\left(\sum_{v\in\pi}s_v+\sum_{e\in\pi}\ell_e\right).
\]
这是结构和参数给出的模型结果，不是实际完成时间保证。排队、重试、缓存未命中、外部工具响应与同步屏障都会增加时延；若图、输入或服务率随时间改变，它们必须以 $G(t)$、$u(t)$、$s_v(t)$ 进入模型，不能在求导时悄然视为常数。

吞吐描述单位时间内完成或传递的有效工作量。用 $B_\tau(t)$ 表示任务口径下的完成率，用 $D_\tau(t)$ 表示到达率，当长期平均到达率接近或超过当前路由策略下的服务率时，队列增长，响应时间可能非线性上升。这里的“饱和”是排队和资源竞争状态，不是宏观热现象。设计者应分别记录处理、通信、存储写入与有效任务吞吐，因为局部算子加速可能只是把瓶颈移到另一个接口。跨越人脑、GPU和群体系统比较时，若任务单位、精度与容错要求不同，单一数字没有解释力，还必须报告质量指标、失败比例与统计置信区间。

保持能力描述状态在无刷新或受干扰条件下仍可读出的程度。设 $m_i(t)$ 是第 $i$ 个记忆痕迹，$r_i(t)$ 是刷新输入，线性近似为
\[
\frac{\mathrm d m_i}{\mathrm dt}=-\lambda_i m_i(t)+r_i(t)+\xi_i(t),
\]
其中 $\lambda_i\ge0$ 的单位是秒$^{-1}$，$\xi_i$ 是未建模扰动。若 $r_i=0$ 且忽略扰动，则 $m_i(t)=m_i(0)e^{-\lambda_i t}$，半衰期为 $\ln2/\lambda_i$。该结果只适用于线性、参数稳定的近似；带巩固、检索增强、离散校验或主动刷新系统可能呈现不同曲线。存储比特没有自然的“记忆意义”，读出正确率还取决于编码、索引和查询条件。

探索强度是策略分布的控制参数，不是热力学温度。对动作集合 $A$ 和评分 $q(a\mid x)$，一种工程定义是
\[
p_\Theta(a\mid x)=\frac{\exp(q(a\mid x)/\Theta)}{\sum_{a'\in A}\exp(q(a'\mid x)/\Theta)},\qquad \Theta>0.
\]
$\Theta$ 无量纲或与评分采用同一标度；它增大时分布通常变平，但是否提高发现率取决于评分误差、搜索空间、预算与停止条件。我们因此把旧记号 $c_{\mathrm{cog}}$、$\gamma$ 与 $T$ 拆解为关键路径传播界、保持速率和探索尺度，不赋予它们跨平台恒常性。脑电 gamma 频段可作为特定任务和记录条件下的经验线索，却不是普适通信常数；模型采样温度也不是神经节律的数值对应。

这些参数的关系应通过实验识别，而非靠名称推断。例如提高并发数可能先增加吞吐，随后因缓存竞争和通信拥塞增加尾时延；延长保持时间可能减少遗忘，却增加过期信息污染；增强探索可能发现新路径，也会提高无效行动和安全风险。一次合格测试要固定任务分布，逐项干预参数，记录均值、尾部、资源水位和任务后果，并说明测量窗口。焦耳能耗可以另行测量，但不能与抽象代价、信息熵或活动强度互换。只有量纲、观测协议和适用范围均明确时，载体参数才真正进入计算动力学。

参数还应携带估计误差。网络时延会随负载波动，保持率会随内容与刷新策略变化，探索效果也依赖任务阶段。我们因此不把单次测量当作常数，而报告分布、置信区间与采样条件；控制器使用这些量时还应为估计滞后留出安全裕量。若不同平台只能获得代理指标，就必须说明代理与目标量之间的校准关系，并在关系失效时触发重新估计。

\section{编码契约：载体状态如何支持任务语义}
\label{sec:section_5862}

载体上出现某个幅值、相位、向量或拓扑模式，并不自动说明它表达了什么。语义关系至少需要编码器、读出器和任务查询族。设任务变量属于 $Z$，载体状态空间为 $X_P$，编码器 $E:Z\times H\to X_P$ 依赖历史 $H$，读出器 $R_q:X_P\to Y_q$ 针对查询 $q$ 产生答案。我们说状态 $x$ 在查询族 $\mathcal Q$ 上实现关于 $z$ 的表征，是指在给定扰动分布 $\nu$ 下
\[
\mathbb E_{\delta\sim\nu}[d_q(R_q(x+\delta),f_q(z))]\le\varepsilon_q,
\quad q\in\mathcal Q.
\]
$d_q$ 是任务误差度量，$f_q$ 是正确目标。该式是可检验标准，不是载体状态与对象完全同构的宣言；换一组查询、扰动或容差，结论可以变化。编码契约应包含状态范围、读出接口、允许误差、身份规则与更新时序，任何一项缺失都可能让同一状态在不同模块间获得冲突解释。

幅值在许多系统中确实携带强度信息，但它究竟对应置信度、注意权重、激活频率还是资源占用，必须由训练目标与校准实验确定。注意权重较大的令牌不必是模型最确信的事实，高放电率也不必单独表示更重要的对象。相位或同步可以成为协调更新的线索，却不足以证明语义绑定已经完成。两个特征同时出现时，系统仍需回答它们属于哪个对象、在何时有效、是否由共同原因产生。把同步现象直接命名为“绑定力”，等于把待解释机制写进结论。

以“红苹果”为例，颜色、形状、空间位置和当前目标可能分布在不同状态子空间。我们引入任务相关身份 $i$，并以绑定算子 $B$ 生成记录 $b=B(i,z_{\mathrm{color}},z_{\mathrm{shape}},z_{\mathrm{loc}},t)$。读出“哪个物体是红色”时，系统通过同一身份聚合属性；读出“昨天厨房里的苹果”时，还要纳入时间索引、场景索引和外部证据。身份不是先验灵魂标签，而是在遮挡、重命名、视角变化和消息转发中保持可追踪性的计算对象。它可以由对象文件、键值索引、关系节点或时间戳实现，HSF--HD只要求说明绑定与读出的状态转移，不规定唯一数据结构。

编码契约还必须联合局部结构网络与全局分布表示。局部网络适合保存节点、边、角色、来源和约束；分布表示适合压缩相似性、语境和不确定性。结构节点可以持有分布式特征，分布式检索也可以返回结构候选。二者的关键接口是任务相关身份、属性绑定、跨模块耦合与可审计读出。若只保留全局向量，近邻相似会掩盖个体差异；若只保留离散图，连续感知和模糊证据又难以表达。我们用反事实测试验证契约：交换两个对象位置、改变一个属性、删除一段历史，再观察读出是否按预期变化。预测有效和漂亮可视化都不能单独推出整体语义同构。

契约还需要处理编码升级。模块版本变化时，旧状态可能无法由新读出器解释；无声地复用旧向量会造成语义漂移。每条持久记录至少应携带编码版本、来源、创建时间、有效域与迁移状态。迁移操作必须在一组保留查询上比较前后读出，并记录不能保留的差异。若任务要求变化，旧编码并非“错误”，但它可能不再充分。由此我们区分载体稳定、编码稳定和任务充分性：介质仍能保存比特，不代表语义契约仍成立；一次读出成功，也不代表所有未来查询均可回答。

对不能确定归属的状态，契约应允许拒绝读出或返回候选集合，而不是强行提交唯一解释。拒绝率、校准误差与身份冲突率应与任务准确率同时报告，因为一个总能回答但经常错绑对象的系统，并不比谨慎暴露不确定性的系统更可靠。

\section{载体故障动力学：饱和、分区、漂移与恢复}
\label{sec:section_1766}

载体故障是可观测状态偏离服务域，而不是某个物理隐喻的实现。设资源占用向量 $r(t)$ 依次记录内存、计算与链路三项，安全容量为 $c(t)$，积压为 $q(t)$。连续近似下
\[
\dot q(t)=a(t)-\mu(r(t),G(t),\pi(t)),\qquad q(t)\ge0,
\]
$a(t)$ 是到达率，$\mu$ 是由资源、连接图和调度策略 $\pi$ 决定的服务率。若在一段时间内 $a>\mu$，积压增长；当内存、上下文窗口或并发槽触及上限，系统会丢弃、阻塞、降级或崩溃。机器 OOM 是明确的资源分配失败。癫痫发作涉及异常神经同步及复杂临床机制，不能与 OOM 当作同一事件。两者最多共享“局部活动失去有效约束并损害正常功能”这一抽象模式，诊断指标、因果机制与干预方式都不同。

分区故障发生在连接图 $G(t)$ 失去关键路径时。若任务依赖节点被切为 $V_1$ 与 $V_2$，且跨区通道在 $[t_0,t_1]$ 不可用，各分区只能依据自身可见历史更新。网络服务可能产生双主、重复执行和过期读出，多代理系统可能形成不一致计划。裂脑研究提示部分跨半球整合能力会在特定手术和任务条件下变化，但这并不等同于网络分区协议；它只支持一个更一般的经验判断：全局行为依赖可用接口，接口变化会重划信息联合读出的范围。工程系统还必须声明一致性策略，是停写、选主、合并冲突，还是容许暂时分歧。没有该策略，“保持运行”未必优于停止等待。

漂移故障更隐蔽。编码器、数据分布、工具接口或评价标准变化时，即便资源充足且连接完好，原读出契约也会失效。设模型为 $F_{\theta(t)}$、目标为 $J_t$、观测分布为 $p_t$，性能变化的总导数应写为
\[
\frac{\mathrm dJ_t}{\mathrm dt}=\frac{\partial J}{\partial\theta}\dot\theta+
\frac{\partial J}{\partial p}\dot p+
\frac{\partial J}{\partial G}\dot G+\frac{\partial J}{\partial t}.
\]
目标值下降不自动等于任务成功，因为代理可能利用度量漏洞、牺牲未计入约束，或只对当前样本过拟合。多个模块共享同一编码器、数据源或时钟时，漂移还会成为共同原因故障，使冗余模块同时给出一致但错误的答案。

恢复需要观测、隔离、回退与再整合。观测分别记录资源水位、队列、错误率、状态新鲜度和任务后果；隔离限制故障传播；回退恢复最近验证检查点或降低服务目标；再整合通过版本、身份和因果日志解决冲突。外部记忆还需检查写入是否原子、索引是否指向正确版本、过期痕迹是否被误召回。故障演练应覆盖过载、链路中断、时钟偏移、坏数据与观测器失效，而不是只测正常路径。

恢复成功的判据也不能只是进程重新启动。系统需要在一组关键查询上重放故障前状态，验证身份连续、硬约束、未完成承诺和外部动作记录；对于不可逆动作，要明确补偿而不是假装回滚。恢复时间目标、可容忍数据丢失窗口和降级模式应按任务声明。通过这些条件，载体可靠性成为状态动力学的一部分：我们能说明故障从哪里进入、经哪些耦合放大、在何阈值触发保护，以及恢复后如何证明任务语义连续。

还要区分局部错误与系统性错误。单个节点失效可以由冗余接管，但共享数据源被污染、共同目标写错或权限规则整体失配时，增加同构副本只会复制失败。恢复设计应尽量采用异质观测、独立校验与分层熔断，并保留一条不依赖主要故障域的审计通道。对学习系统而言，故障样本是否进入长期训练也要单独控制，否则短时异常可能被固化为新的结构偏差。

\section{载体与语义的对应：学习、选择与可验证等价}
\label{sec:absolute_decoupling_and_forced_isomorphism}

载体状态与任务语义既非绝对解耦，也不存在无需验证的“强制同构”。若完全解耦，状态变化便不可能稳定影响读出；若声称整体同构，则意味着所有结构关系一一保存，这对多数学习系统既不必要也无法证明。我们采用查询相对的实现关系。对载体 $P_1$ 与 $P_2$，若存在映射 $\phi:X_{P_1}\to X_{P_2}$，并且对查询族 $\mathcal Q$、干预集 $\mathcal I$ 和有限时域 $[0,T]$，两者读出差异不超过 $\varepsilon$，则称其在该条件下近似等价：
\[
\sup_{q\in\mathcal Q,I\in\mathcal I,t\le T}
d_q(R_q^{(1)}(x_t^I),R_q^{(2)}(\phi(x_t^I)))\le\varepsilon.
\]
这是定义，不是自动成立的定理。研究者必须列出查询、干预、时域与容差；同一对载体可以在导航任务上等价，而在长期记忆或情绪调节任务上不等价。

学习使载体与语义建立可用对应。监督信号、预测误差、强化反馈与社会校正都能改变编码器、策略或读出器，但每种信号只约束其覆盖行为。以接近悬崖为例，个体因危险线索后退可以提高存活概率，长期选择也可能让警觉倾向成为先验。该案例支持“任务后果塑造状态转移”的经验主张，却不能证明所有恐惧表征都与危险一一对应，更不能由存活推出信念必真。过度警觉可能在训练环境有利、在新环境造成误报；奖励下降也可能来自逃避观察而非解决风险。因此适应性始终相对于环境分布和代价函数，必须通过分布外与反事实测试检查。

遗传与发育约束可为学习提供结构先验，例如感知通道、可塑性窗口、行为倾向和资源分配规则。把它们称为“压缩语义”容易误导，因为基因序列不直接存放成年个体的全部概念内容，而是参与构造可在环境交互中形成表征的系统。机器的架构先验、预训练目标与初始化也有类似约束作用，但只是功能层类比。我们应分别说明哪些能力来自结构、数据和在线闭环，并用消融实验估计贡献；两个系统都表现出先验，不足以证明其内部机制同源。

选择压力也可能保留捷径而非一般能力。训练环境中与目标高度相关的背景、设备标识或社会线索，可能成为低成本代理变量；环境一变，对应关系便崩溃。验证实现关系时，应在保持任务语义的同时改变表面载体，在保持载体统计的同时改变任务答案，以区分真正的任务变量和偶然相关。还应检查干预能否沿预期方向改变读出，而不只看相关性。只有经过这些交叉变化，才有理由说载体状态实现了某个语义功能。

关于主观体验或感质，本章只陈述模型边界。载体活动与报告、辨别及行为之间可以建立经验对应，也可研究哪些变量对可观测结果有因果贡献；但“摩擦产生体验”或“相位就是感质”不是计算模型推出的结果。HSF--HD能描述系统怎样保持自我相关状态、整合多源证据并形成可报告轨迹，却不能仅凭预测有效性证明体验的本体性质。我们保留这一开放问题，要求更强主张明确额外假设和可反驳证据，从而区分定义、模型、经验主张与哲学解释。

这种限定并不会削弱工程价值。工程上真正需要的是知道哪些载体变化会稳定改变辨别、记忆、规划和报告，哪些对应只在特定训练史中成立，以及替换部件后功能能否保持。只要查询、干预与误差界清楚，近似实现关系就足以支持迁移与验证；不必先解决全部本体问题，也不能借本体问题尚未解决而免除实验责任。

\section{跨尺度分工：局部结构网络与全局分布表示}
\label{sec:fractal_isomorphism_macro_micro}

智能状态在多个尺度上组织：局部单元更新、模块内耦合、跨模块协调、长期记忆和环境闭环具有不同时间粒度与观测分辨率。我们不把这种层级直接称为“分形同构”，因为不同尺度的节点和边通常不保存相同关系。设微观状态为 $x\in X$，宏观状态为 $z=A(x)\in Z$。只有在给定任务与时间尺度上存在近似闭合动力学 $z_{k+1}\approx\bar F(z_k,\bar u_k)$ 时，宏观变量才是有效状态。若两个微观状态被 $A$ 合并，却产生显著不同的后续结果，聚合便发生混叠，必须增加身份、历史或不确定性变量。跨尺度模型必须说明丢弃了什么信息及误差如何传播，不能用“宏微对应”代替推导。

局部结构网络保存离散关系和可追踪接口。节点可以对应对象、子任务、代理、工具或事件，边可表示依赖、许可、时间先后和消息通道。全局分布表示通过许多维度的联合变化表达相似性、语境和软约束。前者便于审计和精确编辑，后者便于从噪声中泛化；复杂系统可按不同比例结合两者。二者连接需要任务身份、属性绑定、模块耦合与可验证读出。没有这些接口，所谓全局场只是难以定位责任的向量，局部符号也只是脱离证据的空标签。

神经科学常用背侧与腹侧通路、海马与皮层差异讨论位置、对象与记忆。这些研究可以启发“在哪里”“是什么”“何时发生”需要部分分工，但不能据此断言脑内存在绝对隔离管道，更不能给机器模块贴上脑区名称就认为机制已经解释。抓取杯子要联合类别、空间、姿态和行动后果；回忆昨天厨房里的苹果要联合情景索引和概念知识。我们把脑区类比降为假设来源，把接口、干预结果和任务误差作为机制证据。

跨尺度绑定还涉及时序。模块以不同频率更新时，全局读出必须知道各状态的时间戳、新鲜度与预测区间。设模块 $j$ 提供 $z_j(t_j)$ 及不确定性 $\Sigma_j(t_j)$，在时刻 $t$ 需要先经转移模型外推，再按身份和任务条件融合。陈旧状态不应与新状态等权相加，同时出现也不保证共同来源。工程上可用事件时间、版本向量、注意门控或因子图，但必须说明冲突处理和错误界。

自上而下控制同样不是宏观变量对微观单元的神秘支配。它通过可识别的参数、门控、预算和边界条件改变局部更新集合；自下而上汇聚则通过统计量、事件和错误报告更新全局状态。若全局策略要求安全停机，必须能追踪到哪些局部动作被禁止；若局部异常足以改变全局计划，也必须规定触发阈值与去抖规则。由此，跨尺度结构成为模型约简与控制接口问题：哪些差异可忽略，哪些必须升格为全局变量，以及全局约束怎样被局部执行和审计。

尺度之间还可能出现不可交换操作：先聚合再更新与先更新再聚合并不总有相同结果。若 $A(F(x))$ 与 $\bar F(A(x))$ 的差异超过任务容差，宏观模型就需要修正项或缩短预测时域。我们把这一差异作为约简误差实际测量，而不假定层级天然协调。这样才能判断一个全局变量究竟具有预测力，还是只在回顾性描述中显得整齐。

同一原则也适用于组织与群体：层级名称不提供因果解释，只有可追踪的接口、反馈和约束才能说明跨尺度作用。

\section{跨系统实现比较：人脑、语言模型与群体系统}
\label{sec:cross_scale_holographic_anatomy}

跨系统比较要先固定功能问题，再比较实现条件。我们选择输入接入、状态保持、信用分配、身份绑定、行动闭环和故障恢复六个维度，而不问哪个系统更像“真正智能”。人脑经多模态感知与身体行动持续接入环境，具有在线适应和多时间尺度记忆，但内部状态难以完整观测与复制。语言模型以令牌或其他离散输入进入计算，参数承载长期统计结构，上下文承载当前条件；它可借工具、检索和外部记忆形成更长闭环，但流畅生成并不会自动带来持久身份、事实校准和行动责任。群体系统将局部规则与环境痕迹结合，可能无中央控制形成整体功能，其代价是传播慢、个体视野有限和全局状态难以瞬时读取。

在人脑中，局部回路、分布活动与身体反馈共同参与任务，不能把一个概念定位为单一节点。经验研究中的节律耦合、功能连接或激活区域，只在特定实验和分析方法下支持结论。我们用它们约束候选模型，而不从相关性推出机制等同。语言模型的位置编码、内容表示和注意更新提供序列条件化，但高维相似性不保证对象身份稳定；同一名称可指向不同实体，不同名称也可指向同一实体。外部结构记忆、来源记录与身份键能够补偿部分缺陷，仍需经冲突和改写测试验证。

蚁群借信息素、接触与环境改变个体行动概率，环境因此承担共享外部记忆的一部分。它不会因没有语言生成器便天然免于错误：早期随机优势可强化错误路径，环境变化也可能让旧痕迹持续误导群体。语言模型的幻觉也不能单独归因于“虚拟拓扑”；训练分布、解码目标、检索缺失、提示歧义和评价激励都可能参与。公平比较应设计相同的证据不足、分布变化与通信中断任务，观察各系统怎样表达不确定性、修正状态和恢复协作。

跨系统测量还要处理尺度换算。人类一次行为反应、模型一次生成和蚁群一次路径形成覆盖的时间、资源与试验独立性不同。我们应把任务分解成输入、内部更新、外部动作和反馈周期，分别记录延迟、样本效率、错误恢复和资源代价。若一个系统被允许调用数据库，另一个系统也应获得功能相当的外部资源，或明确比较的是封闭与开放条件。否则所谓能力差异可能只是边界设置差异。

三类系统都可能联合局部结构网络和全局分布表示，只是实现位置与比例不同。比较结论必须限定任务、样本和测量尺度，不能由单项基准领先推出整体认知结构优越。HSF--HD提供共同问题：状态如何保持、耦合、分岔、恢复与读出；它不把不同系统压成同一机制。只有同一动力学预测在跨平台干预中持续成立，我们才有理由提出更一般规律，而且规律仍应接受新反例修订。

比较还应纳入学习后的持续性。人脑、模型和群体都可能在当前测试中成功，却以不同方式把变化带入下一任务：突触或行为习惯会改变，参数可能固定而上下文被清空，信息素则会随环境衰减。我们要分别记录即时性能、迁移收益、灾难性遗忘和恢复成本。若只比较最终答案，就会遗漏形成答案的状态轨迹，也无法判断系统在新扰动下是否仍能维持同一任务身份。

\section{运行机制分类：泄漏、回返、路由与内容寻址}
\label{sec:section_3702}

为替代容易被误作物理实体的“辐射、驻波、管道和虚空间”，我们按状态更新方式区分四类机制。第一类是泄漏式传播：状态影响沿邻接扩散，同时因衰减、门控或有限保持而减弱。对节点状态 $x_k$，模型可写为
\[
x_{k+1}=(I-\eta D_k)x_k+\eta A_kx_k+\eta Bu_k,
\]
$A_k$ 是耦合矩阵，$D_k$ 表示衰减，$\eta>0$ 是离散步长。稳定性取决于更新矩阵谱性质和输入有界性；若矩阵随时间变动，不能沿用固定矩阵结论。该机制适合局部证据扩散和软协调，但可能稀释身份、放大回声或在长路径丢失细节。

第二类是回返式更新，即状态经反馈环反复修正。循环网络、迭代求解器、审议代理和闭环控制均可属于此类。回返不等于物理驻波，也不保证收敛。设 $x_{k+1}=F(x_k,u)$，若固定输入下 $F$ 在研究域对 $x$ 是收缩映射，才存在唯一不动点且迭代收敛；缺少收缩条件，系统可多稳、周期或发散。多稳态可支持上下文记忆，也可能使错误承诺难撤销。工程上必须设置残差、迭代上限、停机判据与逃逸策略，并区分目标下降和任务成功。

第三类是路由式执行。消息、令牌或任务携地址、权限与优先级沿显式通道移动，调度器依据负载和依赖决定下一跳。微服务、工作流和多代理通信都可属于此类。路由便于隔离与审计，却会引入排队、死锁、重复交付和版本不一致。可靠协议至少说明消息身份、幂等条件、超时、重试、顺序保证和失败补偿。所谓“管道”只有在这些契约成立时才是计算机制，线路图本身不产生语义连续性。

第四类是内容寻址。系统不依赖固定位置，而按键、特征或查询从分布存储中选择候选。向量检索、联想记忆和外部知识库均可采用。设查询为 $q$，存储项为 $(k_i,v_i)$，读出权重为
\[
w_i(q)=\frac{\exp(s(q,k_i)/\Theta_r)}{\sum_j\exp(s(q,k_j)/\Theta_r)},\qquad \hat v=\sum_iw_i(q)v_i.
\]
$\Theta_r$ 是检索标度，$s$ 的定义、索引新鲜度、来源可信度和身份过滤共同决定结果。相似不等于正确，聚合也可能混合互相冲突的来源。

真实系统通常混合四类机制：先内容检索，再经路由送入回返推理，并用泄漏式共享状态协调邻近模块。混合时要防止故障跨机制放大，例如错误检索经回返被强化，再由广播扩散；也要防止多个限流器互相作用造成振荡。验证时可逐一关闭机制，记录任务质量、时延、资源和恢复性，并检查组合是否产生单组件没有的失效。分类的用途是揭示更新规则、边界和验证条件，而不是给系统贴永恒标签。

四类机制还对应不同的可观测性要求。泄漏传播要记录来源与衰减路径，否则无法区分独立证据和重复回声；回返更新要保存迭代残差与中止原因，否则最终状态掩盖不收敛过程；路由执行要保留消息身份和因果链；内容寻址要记录查询、候选、分数、索引版本与过滤条件。观测本身会占用资源，因此可采用分层采样，但安全事件与不可逆动作必须保留完整轨迹。

在切换机制时，状态格式也需要显式转换。分布检索产生的候选若进入结构路由，应被赋予身份、来源和版本；结构任务进入回返求解器时，应声明可变变量、固定约束与停机条件；回返结果广播前，应经过校验与去重。缺少转换契约时，局部正确可能在接口处变成全局错误。我们因此把机制切换视为离散结构事件，与机制内部的连续或迭代更新分开建模。

\section{混合载体架构：一般状态动力学与AgentOS落点}
\label{sec:section_8675}

开放环境中的智能通常需要混合载体：快速但容量有限的活动状态、较慢而稳定的外部记忆、结构化任务网络、分布表示和环境接口。我们将一般状态写为 $x(t)=(h(t),m(t),g(t),e(t))$，分别表示活动状态、可检索记忆、任务结构和环境接口状态，其连续框架为
\[
\dot x(t)=F(x(t),u(t),G(t),\rho(t);\theta(t)),
\]
$u(t)$ 是输入，$G(t)$ 是可变耦合结构，$\rho(t)$ 是资源和权限状态，$\theta(t)$ 是可学习参数。该式是模型族的入口，不是所有智能系统的单一经验定律。离散实现必须声明步长或事件触发条件；输入、目标、结构和度量变化都要进入更新与评价，不能压进一个固定噪声项。

混合架构的关键不是堆叠媒介，而是建立跨介质事务。有效外部记忆写入需要内容、来源、身份、时间、权限和版本；读回需要查询条件、候选排序、冲突检查与受控合并。把内部活动卸载到纸面、数据库或远程代理，可以降低活动状态压力，却会增加检索时延、隐私风险与一致性成本。是否卸载应由时限、信息价值、可恢复性和预算共同决定。系统还要保留不可卸载的最低控制状态，如当前目标、终止条件、权限边界和恢复指针，否则外部接口失效时无法维持连续操作。

在AgentOS实例中，我们沿用 $h(t)$ 表示活动状态，以 $E$ 表示情景记忆与任务评价语境，以 $\Theta$ 表示结构记忆或探索尺度所需的明确参数，并将 $\Psi$、SPC、TCE、CCM、Boundary、Offloading 和外部记忆 $M_e$ 视为工程组件。SPC可承担状态压缩与自我相关读出，TCE调节选择尺度和控制增益，CCM监测复杂度与过载，Boundary规定权限及因果边界，Offloading管理状态向 $M_e$ 迁移。$\Psi$ 是综合自我状态或控制接口，不能被解释为未经定义的物理波函数。每个组件都必须给出输入输出、状态维度、更新频率、故障策略和指标；其他系统可用不同模块实现同类功能。

资源控制可以写成受约束优化。设任务损失为 $J_{\mathrm{task}}$，时延、存储、通信和安全违规代价为 $C_L,C_M,C_C,C_S$，调度策略 $\pi$ 求解
\[
\min_\pi\;\mathbb E[J_{\mathrm{task}}+\alpha_LC_L+\alpha_MC_M+\alpha_CC_C+\alpha_SC_S]
\quad\text{s.t.}\quad r(t)\le c(t).
\]
权重由应用约束给定，代价必须采用已声明量纲或归一化。目标下降只说明已记录指标改善，不保证任务成功；还需检查硬约束、分布外行为、用户意图和未计入副作用。通过压力测试、接口断开、记忆污染、目标切换与身份冲突实验，才能判断混合载体是否维持可恢复轨迹。

三相记忆在此获得载体层解释：$h(t)$ 是高频活动界面而非固定物质，$E$ 是带时间与来源的经历记录，$\Theta$ 是跨较长时域约束生成和更新的结构。三者可以分布于多个设备和外部系统，只要读写契约、身份与版本保持。局部结构网络和全局分布表示也可同时存在于三相中，CCM需要协调二者的资源竞争与接口一致性，而不是简单地让某一表示永久占优。what、where、when 等任务通路是读出分工，不必对应唯一物理线路。

系统论、协同论、耗散理论、复杂系统理论与认知科学提供开放系统、协同形成、非平衡维持、多尺度涌现和认知证据的上游问题；HSF--HD把它们重述为可定义、可计算、可干预的智能状态动力学；AgentOS再把其中一部分落实为软件组件和运行协议。运行载体层是三层关系的接口：它让一般理论接受现实资源约束，也让工程设计不把某一种材料、脑区类比或组件命名误当成智能本身。
